\documentclass[english,twoside, openright, hidelinks, 11 pt, class=report,crop=false]{standalone}
\input{../bm}
\input{../../preamb}


\begin{document}
\section{Faktorisering}
\reg[\kvadset \label{kvadset}]{
For to reelle tall $ a $ og $ b $ er
\alg{
&&(a+b)^2 &= a^2+2ab+b^2 && \text{(1. kvadratsetning)}\vn
&&(a-b)^2 &= a^2-2ab+b^2 && \text{(2. kvadratsetning)} \vn
&&(a+b)(a-b)&= a^2-b^2 && \text{(3. kvadratsetning)}
}
}
\spr{
$ (a+b)^2 $ og $ (a-b)^2 $ kalles \outl{fullstendige kvadrat}.\vsk	
	
3. kvadratsetning kalles også \outl{konjugatsetningen}.\vsk

Alle kvadratsetningene er \textit{identiteter}. En \outl{identitet} er en ligning som er oppfylt uansett hvilke verdier man gir variablene i ligningen. 
}
\eks[1]{
Skriv om $ a^2+8a+16 $ til et fullstendig kvadrat.

\sv \vsb
\alg{
a^2+8a+16 &= a^2+2\cdot4a+4^2 \\
&= (a+4)^2
}
}
\eks[2]{
Skriv om $ {k^2+6k +7} $ til et uttrykk der $ k $ er et ledd i et fullstendig kvadrat.

\sv \vsb 
\alg{
k^2+6k +7 &= k^2+2\cdot3k+7  \\
&= k^2+2\cdot3k+ 3^2-3^2+7 \\
&= (k+3)^2-2
}
}
\newpage
\eks[3]{
Faktoriser $ x^2-10x+16 $.

\sv
Vi starter med å lage et fullstendig kvadrat:
\alg{
x^2-10x+16 &= x^2-2\cdot5x+5^2-5^2+16 \\
&= (x-5)^2-9
}
Vi legger merke til at $ 9=3^2 $, og bruker 3. kvadratsetning:
\alg{
(x-5)^2-3^2 &= (x-5+3)(x-5-3) \\
&= (x-2)(x-8)
}
Altså er 
\[ x^2-10x+16=(x-2)(x-8) \]
} 
\fork{\ref{kvadset} \kvadset}{
Kvadratsetningene følger direkte av distributiv lov ved multiplikasjon (se \mb).
}
\vsk

\reg[\aenato \label{a1a2}]{
Gitt $ x, b, c \in \mathbb{R} $. Hvis $ a_1+a_2=b $ og $ a_1a_2=c$, er 
\begin{equation}
x^2+bx+c=(x+a_1)(x+a_2)	\label{a1a2eq}
\end{equation}
}
\eks[1]{
Faktoriser uttrykket $ x^2-x-6 $.

\sv
Siden $ 2\cdot(-3)=-6 $ og $ 2+(-3)=-1 $, er
\[ x^2-x-6=(x+2)(x-3) \]
}
\newpage
\eks[2]{
	Faktoriser uttrykket $ b^2-5b+4 $.
	
	\sv
	Siden ${(-4)\cdot( -1)=4} $ og ${ -4+(-1)=-5} $, er 
	\[ b^2-5b+4 =(b-4)(b-1) \] 
}
\fork{\ref{a1a2} \aenato}{
Vi har at
\algv{
(x+a_1)(x+a_2)&=x^2+xa_2+a_1x+a_1a_2 \\
&= x^2+(a_1+a_2)x+a_1a_2
}
Hvis $ a_1+a_2=b $ og $ {a_1a_2=c}$, er
\nn{
(x+a_1)(x+a_2) = x^2+bx+c
}
}
\newpage
\section*{Fortegnsskjema}
En funksjon $f(x)$ består gjerne av uttrykk som har positive og negative verdier for forskjellige verdier av $x$. For å få et tydeligere bilde av når $f$ er positiv eller negativ, kan vi lage et \outl{fortegnsskjema}.
\eks[1]{
	Løs ulikheten
	\[ x^2-8x-9\leq0 \] \vs
	\sv
	Siden $ 1\cdot(-9)=-9 $ og $ 1+(-9)=-8 $, er
	\[ x^2-8x-9=(x+1)(x-9) \]
	Vi setter $ f=(x+1)(x-9) $, og lager et fortegnsskjema:
	\fig{fortgnskj2}
	Skjemaet viser følgende:
	\begin{itemize}
		\item Uttrykket $ x+1 $ er negativt når  $ x<-1 $, lik 0 når $ x=-1 $, og positivt når $ x>-1 $.
		\item Uttrykket $ x-9 $ er negativt når $ x<9 $, lik 0 når $ x=9 $, og positivt når $ x>9 $.
		\item Siden $ f=(x+1)(x-9) $, er \vspace{-5pt}
		\begin{center}
			\begin{tabular}{l}
				$ f>0 $ når $ {{x\in[-\infty, -1)} \cup (9, \infty]} $ \br
				$ f=0 $ når $ {x\in\{-1, 9\}} $ \br
				$ f<0 $ når $ x\in(-1, 9)$
			\end{tabular}
		\end{center}
	\end{itemize}
	Altså er $ x^2-8x-9\leq0 $ når $ x\in[-1, 9] $.
}

\newpage
\eks[2]{
Løs ulikheten
\[ \frac{(x-3)(x+2)}{x-4}>0 \] \vs
\sv
Vi setter $f=\frac{(x-3)(x+2)}{x-4}$, og lager et fortegnskjema:
\fig{fortgnskj3}
Skjemaet viser følgende:
\begin{itemize}
	\item Uttrykket $ x-3 $ er negativt når  $ x<3 $, lik 0 når $ x=3 $, og positivt når $ x>2 $.
	\item Uttrykket $ x+2 $ er negativt når $ x<-2 $, lik 0 når $ x=-2 $, og positivt når $ x>-2 $.
	\item Uttrykket $x-4$ er negativt når $x<4$, lik 0 når $x=4$, og positivt når $x>4$.
	\item Siden $f=\frac{(x-3)(x+1)}{x-4}$, har vi at \vspace{-5pt}
	\begin{center}
		\begin{tabular}{l}
			$ f>0 $ når $ {{x\in(-2, 3)} \cup (4, \infty]} $ \br
			$f$ er ikke definert for $x=4$ \br
			$ f=0 $ når $ {x\in\{-2, 3\}} $ \br
			$ f<0 $ når $ x\in(-\infty, -2)\cup(3, 4)$
		\end{tabular}
	\end{center}
\end{itemize}
Altså er ${\frac{(x-3)(x+2)}{x-4}>0}$ for $ x\in(-2,3)\cup (4, \infty) $.\vsk

\mers{Siden ${x-4}$ kan være både positiv og negativ, kan vi ikke multiplisere ulikheten med denne faktoren (fordi vi ikke kan vite om retningen på ulikhetstegnet endres eller ikke).}
}

\section{Andregradslikninger}
\reg[Andregradslikning med konstantledd]{
	Gitt likningen
	\begin{equation}\label{abc}
		ax^2+bx+c=0
	\end{equation}
	hvor $ a, b, c \in \mathbb{R}$. Da er 
	\begin{flalign*}
		&&\qquad\qquad\qquad	x&= \frac{-b\pm\sqrt{b^2-4ac}}{2a} &&\qquad(\textit{abc}\text{\,-formelen})
	\end{flalign*}
	Hvis $ {x=x_1} $ og $ {x=x_2} $ er løsningene gitt av \textit{abc}-formelen, er
	\begin{equation}\label{abcfakt}
		ax^2+bx+c=a(x-x_1)(x-x_2)
	\end{equation}
}
\eks[1]{
	\abc{
	\item Løs likningen	$ 2x^2-7x+5=0  $.
	\item Faktoriser uttrykket på venstre side i oppgave a). 
} \vs
\sv \vs
\abc{
\item Vi bruker \textit{abc}-formelen. Da er $ {a=2} $, $ {b=-7} $ og $ {c=5} $. Dermed er
\alg{
	x &=\frac{-(-7)\pm\sqrt{(-7)^2-4\cdot2\cdot5}}{2\cdot2} \br
	&=\frac{7\pm\sqrt{49-40}}{4} \br
	&=\frac{7\pm\sqrt{9}}{4}\br
	&=\frac{7\pm3}{4}	
}
Enten er
\[ x=\frac{7+3}{4}=\frac{5}{2} \]	
eller så er
\[ x=\frac{7-3}{4}=1 \]	

\item $ 2x^2-7x+5=2(x-1)\left(x-\frac{5}{2}\right) $ 	
}
}
\eks[2]{
	Løs likningen 
	\[ x^2+3x-10=0 \]
	\sv
	Vi bruker \textit{abc}-formelen. Da er $ a=1 $, $ b=3 $ og $ c=-10 $. Dermed har vi at
	\alg{
		x&=\frac{-3\pm\sqrt{3^2-4\cdot1\cdot(-10)}}{2\cdot1}\br
		&=\frac{-3\pm\sqrt{9+40}}{2} \br
		&=\frac{-3\pm7}{2}
	}
	Altså er
	\[ x=-5\qquad\vee\qquad x=2 \]
}
\eks[3]{
Løs likningen	
\[ 4x^2-8x+1=0 \]
\sv
Av \textit{abc}-formelen har vi at
\alg{
x&=\frac{8\pm\sqrt{(-8)^2-4\cdot4\cdot1}}{2\cdot4} \\
&=\frac{8\pm 4\sqrt{4-1}}{8} \br
&=\frac{2\pm\sqrt{3}}{2}
}
Altså er 
\[ x=\frac{2+\sqrt{3}}{2}\quad\vee\quad \frac{2-\sqrt{3}}{2} \]
}

\newpage
\fork{Andregradslikninger}{
	Gitt likningen
	\[ ax^2+bx+c=0 \]
	Vi starter med å omskrive likningen:
	\[ x^2+\frac{b}{a}x+\frac{c}{a}=0 \]
	Så lager vi et fullstendig kvadrat, og anvdender konjugatsetningen til å faktorisere uttrykket:
	\alg{
		x^2+\frac{b}{a}x+\frac{c}{a}&=x^2+2\cdot\frac{b}{2a}x+\frac{c}{a}\\
		&=\left(x+\frac{b}{2a}\right)^2-\frac{b^2}{4a^2}+\frac{c}{a} \\
		&=\left(x+\frac{b}{2a}\right)^2-\frac{b^2-4ac}{4a^2} \\
		&=\left(x+\frac{b}{2a}\right)^2-\left(\sqrt{\frac{b^2-4ac}{4a^2}}\,\right)^2 \br
		&=\left(x+\frac{b}{2a}\right)^2-\left(\frac{\sqrt{b^2-4ac}}{2a}\,\right)^2 \br
		&=\left(x+\frac{b}{2a}+\frac{\sqrt{b^2-4ac}}{2a}\,\right)\left(x+\frac{b}{2a}-\frac{\sqrt{b^2-4ac}}{2a}\,\right)
	}	
	Uttrykket over er lik 0 når
	\[ x=\frac{-b+\sqrt{b^2-4ac}}{2a}\qquad\vee\qquad x=\frac{-b-\sqrt{b^2-4ac}}{2a} \]
}
\newpage
\section{Polynomdivisjon}
\subsection*{Metoder}
Når to tall ikke er delelige med hverandre, kan vi kombinere et heltall og en brøk for å uttrykke kvotienten. For eksempel er
\begin{equation}\label{17div3}
	\frac{17}{3}=5+\frac{2}{3}
\end{equation}
Tanken bak \eqref{17div3} er at vi skriver om telleren slik at vi får fram den\\ delen av 17 som er delelig med 3:
\[ \frac{17}{3}=\frac{5\cdot3+2}{3}=5+\frac{2}{3} 
\]
Den samme tankegangen kan brukes for brøker med polynomer, og da kalles det \outl{polynomdivisjon}. 
\regv
\eks[1]{ \label{polydiveks1}
Utfør polynomdivisjon på uttrykket
\[ \frac{2x^2+3x-4}{x+5} \] 
\sv
\metode{Metode 1}{0.4\linewidth} \\
Vi gjør følgende trinnvis; med den største potensen av $ x $ i telleren som utgangspunkt, lager vi uttrykk som er delelige med telleren.
\alg{
	\frac{2x^2+3x-4}{x+5}&=\frac{2x(x+5)-10x+3x-4}{x+5} \br
	&= 2x+\frac{-7x-4}{x+5} \br
	&= 2x+\frac{-7(x+5)+35-4}{x+5} \br
	&= 2x-7 +\frac{31}{x+5}
}
\newpage
\metode{Metode 2}{0.4\linewidth} \\
(Se utregningen under punktene)
\begin{enumerate}[label=\roman*)]
	\item  Leddet med den høyste graden av $ x $ i dividenden er $ 2x^2 $. Dette uttrykket kan vi få fram ved å multiplisere dividenden med $ 2x $. Vi skriver $ 2x $ til høyre for likhetstegnet, og subtraherer $ 2x(x+5)=2x^2+10x $.
	\item Differansen fra punkt i) er ${-7x-4}$. Vi kan få fram leddet med den høyeste gradenen av $ x $ ved å multiplisere dividenden med $ -7 $. Vi skriver $ -7 $ til høyre for likhetstegnet, og subtraherer $ -7(x+5)=-7x-5 $. 
	\item Differansen fra punkt ii) er $ 31 $. Dette er et uttrykk som har lavere grad av $ x $ enn dividenden, og dermed skriver vi $ \frac{31}{x+5} $ til høyre for likhetstegnet.
\end{enumerate}
\alg{
	\phantom{-}&(2x^2+3x-4):(x+5) =2x-7+\frac{31}{x+5} \\ 
	-&\underline{(2x^2+10x)} \\
	&\phantom{02x-}-7x-4  \\
	&\phantom{xx}-\underline{(-7x-35)}\\
	&\phantom{aaaaaaaaaa''\,}31 
}
}
\newpage
\eks[2]{
Utfør polynomdivisjon på uttrykket
\[ \frac{x^3-4x^2+9}{x^2-2} \]
\sv
\metode{Metode 1}{0.4\linewidth}\\
\alg{
\frac{x^3-4x^2+9}{x^2-2}&= \frac{x(x^2-2)+2x-4x^2+9}{x^2-2} \br
&=x+\frac{-4x^2+2x+9}{x^2-2}\br
&=x+\frac{-4(x^2-2)-8+2x+9}{x^2-2}\br
&= x-4+\frac{2x+1}{x^2-2}
}
\metode{Metode 2}{0.4\linewidth} \\
\alg{
	\phantom{-}&(x^3-4x^2+9):(x^2-2) =x-4+\frac{2x+1}{x^2-2} \\ 
	-&\underline{(x^3-2x)} \\
	&\phantom{02x}-4x^2+2x+9  \\
	&\phantom{'}-\underline{(-\,4x^2+8)}\\
	&\phantom{aaaaaaaaaa'}2x+1 
}
}\newpage

\eks[3]{ \label{polydiveks3}
Utfør polynomdivisjon på uttrykket
\[ \frac{x^3-3x^2-6x+8 }{x-4}\]

\sv
\metode{Metode 1}{0.4\linewidth} \\
\alg{
\frac{x^3-3x^2-6x+8 }{x-4} &= \frac{x^2(x-4)+4x^2-3x^2-6x+8}{x-4} \br
&= x^2 +\frac{x^2-6x+8}{x-4} \br
&= x^2+\frac{x(x-4)+4x-6x+8}{x-4}\br
&=x^2+x+\frac{-2x+8}{x-4} \br
&= x^2+x-2
}
\metode{Metode 2}{0.4\linewidth} \\
\alg{
	\phantom{-}&(x^3-3x^2-6x+8):(x-4) =x^2+x-2 \\ 
	-&\underline{(x^3-4x^2)} \\
	&\phantom{02xx''',}x^2-6x+8  \\
	&\phantom{''}-\underline{(-\,x^2-4x)}\\
	&\phantom{aaaaaa''''}-2x+8 \\
	&\phantom{aaaa,,,}-\underline{(-2x+8)}\\
	&\phantom{aaaaaaaaaaaaaa,}0
}
%\mer Dette betyr at
%\[ x^3-3x^2-6x+8 = (x^2+x-2)(x-4) \]
}
\newpage
\subsection*{Delelighet og faktorer}
Eksemplene på side \pageref{polydiveks1}\,-\,\pageref{polydiveks3} peker på noen viktige sammenhenger som gjelder for generelle tilfeller:\regv
\reg[Polynomdivisjon \label{polydiv}]{
La $ A_k $ betegne et polynom $ A $ med grad $ k $. Gitt heltallene $m$ og $n$, hvor $m\geq n>0$ og polynomet $ P_m $. Da finnes polynomene $ Q_n $, $ S_{m-n} $ og $ R_{n-1} $ slik at
\begin{equation}\label{polydiveq}
\frac{P_m}{Q_n}=S_{m-n}+\frac{R_{n-1}}{Q_n}	
\end{equation}
}
\spr{
Hvis $ R_{n-1}=0 $, sier vi at $ P_m $ er \outl{delelig} med $ Q_n $.
} 
\eks[1]{
Undersøk om polynomene er delelige med $ {x-3} $.
\abc{
\item $P(x)= x^3+5x^2-22x-56 $
\item $K(x)= x^3+6x^2-13x-42 $
}
\sv \vs
\abc{
\item Ved polynomdivisjon finner vi at
\[ \frac{P}{x-3}=x^2+8x+2-\frac{50}{x-3} \]
Altså er $ P $ \textsl{ikke} delelig med $ x-3 $.
\item Ved polynomdivisjon finner vi at
\[ \frac{K}{x-3}=x^2+9x+14 \]
Altså er $ K $ delelig med $ x-3 $.
}
}


\newpage
\reg[Faktorer i polynomer \label{polyfakt}]{
	Gitt et polynom $ P(x) $ og en konstant $ a $. Da har vi at
	\begin{equation} \label{divisiblep0}
		P \text{ er delelig med } x-a  \Longleftrightarrow P(a)=0 
	\end{equation}
	Hvis ett av vilkårene over er oppfylt, finnes det et polynom $ S(x) $ slik at
	\begin{equation}\label{PfaktS}
		P=(a-x)S
	\end{equation}
}
\eks[1]{
Gitt polynomet
\[ P(x)= x^3-3x^2-6x+8\]
\abc{
\item Vis at $ x=1 $ løser likningen $ P=0 $.
\item Faktoriser $ P $. 
}
\sv

\abc{
\item Vi undersøker $ P(1) $:
\[ P(1)=1^3-3\cdot1^3-6\cdot1+8 = 0 \]
Altså er $ P=0 $ når $ x=1 $.
\item Siden $ P(1)=0 $, er $ x-1 $ en faktor i $ P $. Ved polynomdivisjon finner vi at
\[ P=(x-1)(x^2-2x-8) \]
Da $ 2\cdot(-4)=-8 $ og $ -4+2=-2 $, er
\[ x^2-2x-8=(x+2)(x-4) \]
Dette betyr at
\[ P=(x-1)(x+2)(x-4) \]
}
}
\section{Eulers tall}
\outl{Eulers tall}\index{eulers tall}\index{e} er en konstant som har så stor betydning i matematikk at den har fått sin egen bokstav \sym{$ e $}. Tallet er irrasjonalt, og de ti første sifrene er
\st{\[ e=2.718281828... \]}
De mest fascinerende egenskapene til dette tallet kommer til syne når man undersøker funksjonen ${f(x)= e^x} $. Dette er en eksponentialfunksjon\footnote{Se kapittel ??} som er så viktig at den rett og slett går under navnet \\\outl{eksponentialfunksjonen}. Denne funksjonen skal vi undersøke nærmere i \refved{eulerstallfork} og \refkap{Derivasjon}.
\fig{ekspfunk}
\section{Logaritmer}
I \mb\, så vi at potenstall består av et grunntall og en eksponent. En \outl{logaritme} er en matematisk operasjon relativ til et tall. Hvis en logaritme er relativ til grunntallet til en potens, vil operasjonen resultere i ekspontenten.\vsk

Logaritmen relativ til 10 skrives \sym{$ \log_{10} $}. Da er for eksempel
\[ \log_{10} 10^2 = 2 \]
Videre er for eksempel
\[ \log_{10} 1000= \log_{10} 10^3 = 3\]
Følelig kan vi skrive
\[ 1000 = 10^{\log_{10} 1000} \]
Med potensreglene som utgangspunkt (se \mb), kan man utlede mange regler for logartimer.\regv
\regdef[Logaritmer]{
La $ \log_a $ betegne logaritmen relativ til $ {a>0}$. For $ m\in\mathbb{R} $ er da
\begin{equation}
	\log_a a^m = m
\end{equation}
Alternativt kan vi skrive \vs
\begin{equation}
	m=a^{\log_{a} m}
\end{equation}
}
\eks[1]{ \vs
	\[ \log_5 5^9 = 9 \]
}
\eks[2]{ \vs
	\[ 3 = 8^{\log_8 3} \]
}
\spr{
	\sym{$ \log_{10}$} skrives ofte som \sym{$ \log $}, mens \sym{$ \log_e $} skrives ofte som \sym{$ \ln $} eller (!) \sym{log}. Når man bruker digitale hjelpemidler til å finne verdier til logaritmer er det derfor viktig å sjekke hva som er grunntallet. I denne boka skal vi skrive \sym{$\log_{10}$} som \sym{$\log$}, og \sym{$ \log_e $} som \sym{$ \ln $}.	\\[5pt]
	
	Logaritmen med $ e $ som grunntall kalles den \outl{naturlige logaritmen}.
}
 \regv
\eks[3]{ \vs
\[ \log 10^7 = 7 \]
}
\eks[4]{ \vs
	\[ \ln e^{-3} = -3 \]
} \vsk

\reg[Logaritmeregler]{
\notesm{Logaritmereglene er her gitt ved den naturlige logaritmen. De samme regelene vil gjelde ved å erstatte \sym{$ \ln $} med \sym{$ \log_a $}, og $ e $ med $ a $ for $ a>0 $.
} \vsk

For $ {x, y>0} $ har vi at \vs
\begin{align}
	\ln e&= 1\label{loga}\vn	
	\ln 1 &= 0 \label{log1}\vn
	\ln (xy)&=\ln x + \ln y \label{logxy} \vn
	\ln \left(\frac{x}{y}\right)&= \ln x - \ln y \label{logxdivy} 
\end{align}
For et tall $ y $ og $ x>0 $, er
\begin{equation}
		\ln x^y = y \ln x \label{logxexpy}
\end{equation}
}
\eks[1]{
\[ \ln\left(ex^5\right) =\ln e+\ln x^5=1+5\ln x\]
}
\eks[2]{
	\[ \ln \frac{1}{2} = \ln 1 - \ln 2 = -\ln 2\]
}
\newpage
\fork{Logaritmerregler}{
	\textbf{Likning} \textbf{(\ref{loga})}		
	\[ \ln e= \ln e^1=1  \]
	
	\textbf{Likning} \textbf{(\ref{log1})}
	\[ \ln 1 = \ln e^0 =0  \]
	
	\textbf{Likning} \textbf{(\ref{logxy})}\\
	For $ m, n\in \mathbb{R} $, har vi at
	\[ \ln e^{m+n}=m+n = \ln e^m + \ln e^n \]
	Vi setter\footnote{Vi tar det her for gitt at alle positive tall forskjellige fra 0 kan uttrykkes som et potenstall.} $ {x=e^m} $ og $ {y=e^n} $. Siden $ {\ln e^{m+n}=\ln (e^m\cdot e^n)}$, er da
	\[ \ln(x y)=\ln x+\ln y \]
	
	\textbf{Likning} \textbf{(\ref{logxdivy})}\\
	Ved å undersøke $ \ln e^{m-n} $, og ved å sette $ {y=e^{-n}} $, blir for-klaringen tilsvarende den gitt for likning \eqref{logxy}.\vsk
	
	\textbf{Likning} \textbf{(\ref{logxexpy})} \os
	Siden $ x=e^{\ln x} $ og\footnote{Se potensregler i \mb.} $ \left(e^{\ln x}\right)^y = e^{y \ln x} $, har vi at
	\[ \ln x^y = \ln e^{y\ln x} = y\ln x \]
}
\newpage
\section{Forklaringer}
\fork{\ref{polydiv} Polynomdivisjon}{
Gitt heltallene $m$ og $n$, hvor $m>n>0$, og polynomene
\alg{
&P_m\text{, hvor }ax^m\text{ er leddet med høyest grad}\vn
&Q_{n}\text{, hvor } bx^n\text{ er leddet med høyest grad}	
}
Da kan vi skrive
\begin{equation*}
	P_m=\frac{a}{b}x^{m-n}Q_n-\frac{a}{b}x^{m-n}Q_n+P_m
\end{equation*}
Vi setter $ {U=-\frac{a}{b}x^{m-n}Q_n+P_m} $, og merker oss at $ U $ \\nødvendigvis har grad lavere eller lik $ m-1 $. Videre har vi at
\begin{equation} \label{polydivu}
\frac{P_m}{Q_n}=\frac{a}{b}x^{m-n}+\frac{U}{Q_n}
\end{equation}
La oss kalle det første og det andre leddet på høyresiden i \eqref{polydivu} for henholdsvis et ''potensledd' og en ''resterende brøk''. Ved å følge prosedyren som ledet oss til \eqref{polydivu}, kan vi også uttrykke $ \frac{U}{Q_n} $ ved et ''potensledd'' og et ''resterende ledd''. Dette ''potensleddet'' vil ha grad mindre eller lik $ {m-1} $, mens telleren i det ''resterende leddet'' vil ha grad mindre eller lik $ {m-2} $. Ved å anvende \eqref{polydivu} kan vi stadig lage nye ''potensledd'' og ''resterende ledd'' fram til vi har et ''resterende ledd'' med grad $ n-1 $.
}
\newpage
\fork{\ref{polyfakt} Faktorisering av polynom}{
\textbf{\boldmath $ P $ er delelig med $ x-a \Rightarrow P(a)=0 $}\os
Gitt at $P$ er delelig med ${x-a}$. For et polynom $ S $ har vi da av \eqref{polydiveq} at
\algv{
\frac{P}{x-a}&= S \\
P &= (x-a)S
}
Da er åpenbart $ {x=a} $ en løsning for likningen $ {P=0} $, og da er $P(a)=0$.\vsk

\textbf{\boldmath $P(a)=0 \Rightarrow$ $ P $ er delelig med $ x-a $}\os
Av \eqref{polydiveq} finnes et polynom $ S $ og en konstant $ R $ slik at
\alg{
\frac{P}{x-a}&=S+\frac{R}{x-a} \br
P &= (x-a)S+R
}
Hvis $ P(a)=0 $, er $R=0$, og da er $ P $ delelig med $ x-a $.
}
\end{document}